% ID: SEL-D0023-OB-E04
% ESTADO: cotejada. Segunda pasada visual completa y borrador compilado revisados.
% Fuente: Selectividad/matematicas_aplicadas/reserva_b_EXAMEN 2 -ANDALUCÍA.pdf
% SHA-256: 8cdb86259405652622f2ea129bffbdb239076ac612fd7fcebb86b9d5bdc451ec
% Página: 2; opción B; ejercicio original 4.

El peso de los paquetes de azúcar de una marca, medido en gramos, sigue una distribución Normal con desviación típica de 16 gramos. A partir de una muestra de 100 paquetes de azúcar de dicha marca, se obtuvo un peso medio de 247 gramos.
\begin{enumerate}[label=\alph*)]
\item \textbf{(1.5 puntos)} Obtenga un intervalo de confianza para el peso medio de los paquetes de azúcar de esa marca, con un nivel de confianza del 97\%.
\item \textbf{(1 punto)} Determine el tamaño muestral mínimo necesario para estimar el peso medio con un error máximo de 0.5 gramos, a un nivel de confianza del 95\%.
\end{enumerate}
